\section*{الفصل \ref{ch:g5:water-shapes-the-land} --- الماء ينحت الأرض}

\begin{solution}{exo:g5:water-shapes-the-land:1}
(أ) \omterm{def:g5:water-shapes-the-land:weathering}{تجوية}: يتكسر \omterm{def:g3:stones-pebbles-sand:rock}{الصخر} في مكانه. (ب) \omterm{def:g5:water-shapes-the-land:erosion}{تعرية}: يُبرى \omterm{def:g3:stones-pebbles-sand:sand}{الرمل}
ويُحمل بعيدًا. (ج) \omterm{def:g5:water-shapes-the-land:deposition}{ترسيب}: يستقر \omterm{def:g3:stones-pebbles-sand:sand}{الوحل}. (د)
\omterm{def:g5:water-shapes-the-land:weathering}{تجوية}: الجذر يكسر الحجر في مكانه.
\end{solution}

\begin{solution}{exo:g5:water-shapes-the-land:2}
\omterm{def:g5:water-shapes-the-land:source-mouth}{المنبع} هو المكان الذي يبدأ منه \omterm{def:g1:landscapes-around-us:water}{النهر} (عين ماء، أو \omterm{def:g1:landscapes-around-us:water}{نهر} جليدي يذوب، أو أرض
عالية رطبة). \omterm{def:g5:water-shapes-the-land:source-mouth}{والمصب} هو المكان الذي ينتهي فيه، في \omterm{def:g2:water-on-earth:ocean}{البحر} أو في \omterm{def:g1:landscapes-around-us:water}{بحيرة} أو في \omterm{def:g1:landscapes-around-us:water}{نهر}
أكبر. \omterm{def:g5:water-shapes-the-land:source-mouth}{والرافد} \omterm{def:g1:landscapes-around-us:water}{نهر} أصغر يصب
في \omterm{def:g1:landscapes-around-us:water}{نهر} أكبر.
\end{solution}

\begin{solution}{exo:g5:water-shapes-the-land:3}
\omterm{def:g5:water-shapes-the-land:source-mouth}{المنبع}؛ \omterm{def:g1:landscapes-around-us:landform}{الوادي} الضيق على شكل حرف V؛ \omterm{def:g5:water-shapes-the-land:meander}{المنعطفات}؛ \omterm{def:g5:water-shapes-the-land:delta}{الدلتا}.
\end{solution}

\begin{solution}{exo:g5:water-shapes-the-land:4}
(أ) خطأ: الصقيع والماء والجذور والحر والبرد تفتت ببطء حتى \omterm{def:g3:stones-pebbles-sand:rock}{الصخر} الصلب
(\omterm{def:g5:water-shapes-the-land:weathering}{التجوية}). (ب) صواب: كلما كان الماء أسرع، كانت القطع التي يستطيع تحريكها
أكبر. (ج) خطأ: أثقل القطع تستقر أولًا؛ \omterm{def:g3:stones-pebbles-sand:sand}{والوحل} الأنعم يستقر
آخرًا، حيث يكاد الماء يسكن، في \omterm{def:g1:landscapes-around-us:water}{بحيرة} أو في \omterm{def:g2:water-on-earth:ocean}{البحر}.
\end{solution}

\begin{solution}{exo:g5:water-shapes-the-land:5}
A في الجهة الخارجية للانعطاف: الماء هناك سريع عميق، فيبري
الضفة (\omterm{def:g5:water-shapes-the-land:erosion}{تعرية}). وB في الجهة الداخلية للانعطاف: الماء
هناك بطيء، فيضع \omterm{def:g3:stones-pebbles-sand:sand}{الرمل} (\omterm{def:g5:water-shapes-the-land:deposition}{ترسيب}).
\end{solution}

\begin{solution}{exo:g5:water-shapes-the-land:6}
في \omterm{def:g1:landscapes-around-us:landform}{الجبال} يكون \omterm{def:g1:landscapes-around-us:water}{النهر} شديد الانحدار سريعًا: فيستطيع أن يدحرج الحجارة الكبيرة،
ويحمل الحبات الصغيرة بعيدًا. وقرب \omterm{def:g2:water-on-earth:ocean}{البحر} يكون \omterm{def:g1:landscapes-around-us:water}{النهر} بطيئًا:
فلا يعود يستطيع تحريك الحجارة الكبيرة، بل \omterm{def:g3:stones-pebbles-sand:sand}{الرمل} \omterm{def:g3:stones-pebbles-sand:sand}{والوحل} وحدهما، فيضعهما
هناك. وفي الطريق أيضًا، تكون الحجارة قد تكسرت وبُريت حتى صارت حبات
أصغر.
\end{solution}

\begin{solution}{exo:g5:water-shapes-the-land:7}
(أ) المجرى الأعلى: $600 \div 6 = 100$ كيلومتر. والمجرى الأوسط:
$600 \div 2 = 300$ كيلومتر. (ب) المجرى الأدنى: $600 - 100 - 300 = 200$
كيلومتر. وبما أن $3 \times 200 = 600$، فهو $\frac{1}{3}$ \omterm{def:g1:landscapes-around-us:water}{النهر}.
(تحقق: $\frac{1}{6} + \frac{1}{2} + \frac{1}{3} = \frac{1}{6} +
\frac{3}{6} + \frac{2}{6} = \frac{6}{6}$.)
\end{solution}

\begin{solution}{exo:g5:water-shapes-the-land:8}
البقع الفاتحة \omterm{def:g3:stones-pebbles-sand:sand}{رمل} \omterm{def:g3:stones-pebbles-sand:pebble}{وحصباء} وُضعا في الجهة الداخلية للانعطافات،
حيث الماء بطيء (\omterm{def:g5:water-shapes-the-land:deposition}{ترسيب}). \omterm{def:g1:landscapes-around-us:water}{والبحيرات} الداكنة المقوسة حلقات قديمة
من \omterm{def:g1:landscapes-around-us:water}{النهر}: اقتربت حلقتان إلى حد أن \omterm{def:g1:landscapes-around-us:water}{النهر} اخترق عنق
الأرض بينهما، وسلك الطريق المختصر، وترك الحلقة القديمة وراءه
بحيرةً.
\end{solution}

\begin{solution}{exo:g5:water-shapes-the-land:9}
المروج المستوية بجانب \omterm{def:g1:landscapes-around-us:water}{النهر} هي \omterm{def:g5:water-shapes-the-land:flood}{سهله الفيضي}: الأرض التي يغطيها \omterm{def:g1:landscapes-around-us:water}{النهر}
حين يفيض. وهي مستوية تحديدًا لأن \omterm{def:g1:landscapes-around-us:water}{النهر} بناها
من وحله. وسيفيض \omterm{def:g1:landscapes-around-us:water}{النهر} يومًا ما من جديد فيملأ المدرسة
بالماء، ومياه \omterm{def:g5:water-shapes-the-land:flood}{الفيضان} السريعة خطرة. فيجب أن تُبنى المدرسة
على أرض أعلى، وأن تبقى المروج حقولًا أو مروجًا أو حديقة.
\end{solution}

\begin{solution}{exo:g5:water-shapes-the-land:10}
(أ) في السنة 12 شهرًا: $12$ مليار $\div\, 12 = 1$ مليار طن في الشهر
في المتوسط. (ب) المليار هو $1\,000\,000\,000$؛ وبالقسمة على 10، يكون
$100\,000\,000$: أي مئة مليون حمولة شاحنة.
\end{solution}

\begin{solution}{exo:g5:water-shapes-the-land:11}
قصة ممكنة. الحبة جزء من صخرة عاليًا على \omterm{def:g1:landscapes-around-us:landform}{جبل}. يفتح الصقيع
شقًا حولها، ليلة بعد ليلة، حتى تنفصل قطعة من
\omterm{def:g3:stones-pebbles-sand:rock}{الصخر}: تلك \omterm{def:g5:water-shapes-the-land:weathering}{تجوية}. وتسقط القطعة في الجدول تحتها. وفي المجرى
الأعلى يكون الماء سريعًا شديد الانحدار: فيدحرج القطعة (\omterm{def:g5:water-shapes-the-land:erosion}{تعرية})، وباصطدامها
بحجارة أخرى تتكسر \omterm{def:g3:stones-pebbles-sand:pebble}{حصى} ثم حبات
\omterm{def:g3:stones-pebbles-sand:sand}{رمل}. وفي المجرى الأوسط يتباطأ \omterm{def:g1:landscapes-around-us:water}{النهر}؛ فتستريح الحبة مدة على
الشاطئ الرملي في الجهة الداخلية لانعطاف، حتى يحملها \omterm{def:g5:water-shapes-the-land:flood}{فيضان} من جديد. وفي
المجرى الأدنى يكون \omterm{def:g1:landscapes-around-us:water}{النهر} بطيئًا عريضًا؛ فتسافر الحبة مع
الماء إلى \omterm{def:g5:water-shapes-the-land:source-mouth}{المصب}. وهناك يكاد الماء يتوقف، فتستقر الحبة
مع \omterm{def:g3:stones-pebbles-sand:sand}{وحل} \omterm{def:g5:water-shapes-the-land:delta}{الدلتا} ورملها: ذلك \omterm{def:g5:water-shapes-the-land:deposition}{ترسيب}.
\end{solution}
